% Clase 14 --- Carga del condensador (§1.4).
% Movida acá desde el cuerpo de notes.tex (ver clase14-circuito-rc).
% Las dos curvas normalizadas, para que se vea que arrancan en extremos
% opuestos y comparten la misma constante de tiempo.
\begin{center}
\begin{tikzpicture}
\begin{axis}[fisfig,
  width=11cm, height=6cm,
  xlabel={$t/RC$}, ylabel={valor normalizado},
  xmin=0, xmax=5, ymin=0, ymax=1.05,
  domain=0:5, samples=100,
  legend pos=north east,
]
  \addplot[figblue] {1 - exp(-x)};
  \addlegendentry{$Q/\varepsilon C$}
  \addplot[figred] {exp(-x)};
  \addlegendentry{$I/(\varepsilon/R)$}
  \draw[figaux] (axis cs:1,0) -- (axis cs:1,1);
  \node[figgray, anchor=south west, font=\footnotesize]
        at (axis cs:1,0.03) {$\tau=RC$};
\end{axis}
\end{tikzpicture}\\[0.5em]
{\small\itshape Las dos curvas son la misma exponencial vista al derecho y al
revés, y eso no es casualidad: $I$ es la derivada de $Q$, y la derivada de una
exponencial es otra exponencial con la misma constante de tiempo. Los valores
iniciales salen de razonar el circuito sin resolver nada: en $t=0$ el
condensador está descargado, no cae potencial en él, se comporta como un cable y
toda la FEM cae en la resistencia, de donde $I(0)=\varepsilon/R$; al final está
totalmente cargado, toda la FEM cae en él y no circula corriente. A $t=\tau=RC$
la carga alcanzó el $63\%$ de su valor final.}
\end{center}
